\section{Methodology}
\label{sec:method}

In this section, we present the formal mathematical derivation of Curvature-Aware Analytical Model Merging (ACM). We begin by formulating the multi-task model merging problem, then introduce our subspace-restricted quadratic approximation of the loss landscape, derive the exact closed-form optimal solution, and finally present an efficient finite-difference estimation scheme for the projected Hessian products.

\subsection{Problem Setting and Formulation}
Let $W_0 \in \mathbb{R}^D$ represent the parameters of a pretrained base neural network. Suppose we are given $K$ task-specific expert models, where each expert $W_k \in \mathbb{R}^D$ (for $k \in \{1, \dots, K\}$) has been independently fine-tuned from the same initialization $W_0$ on its respective task dataset $\mathcal{D}_k$. The task vector for expert $k$ at layer $l \in \{1, \dots, L\}$ is defined as:
\begin{equation}
    v_k^l = W_k^l - W_0^l
\end{equation}
where $W_k^l$ and $W_0^l$ denote the parameters of expert $k$ and the base model at layer $l$, respectively.

In layer-wise model merging, our goal is to find a set of merging coefficients $\Lambda = \{\Lambda^l\}_{l=1}^L$, where $\Lambda^l = [\Lambda_1^l, \dots, \Lambda_K^l]^T \in \mathbb{R}^K$ is the coefficient vector for layer $l$. The merged model weights at layer $l$ are given by:
\begin{equation}
    W^l(\Lambda^l) = W_0^l + \sum_{k=1}^K \Lambda_k^l v_k^l
\end{equation}
We wish to determine the optimal coefficients $\Lambda$ that minimize the joint multi-task supervised loss across all $K$ tasks:
\begin{equation}
    \min_{\Lambda} \mathcal{L}_{\text{joint}}(W(\Lambda)) = \min_{\Lambda} \sum_{k=1}^K \alpha_k \mathcal{L}_k(W(\Lambda))
\end{equation}
where $\mathcal{L}_k$ is the supervised loss of task $k$ on dataset $\mathcal{D}_k$, and $\alpha_k > 0$ represents the task weight (typically $\alpha_k = 1/K$).

\subsection{Second-Order Taylor Approximation and Subspace Projection}
In realistic settings, the fine-tuned task experts $W_k$ do not reach a perfect zero-gradient minimum. To ensure mathematical consistency under non-zero residual expert gradients, we approximate each task loss $\mathcal{L}_k(W(\Lambda))$ around its expert state $W_k$ using a complete second-order Taylor expansion including the first-order linear term:
\begin{equation}
\begin{split}
    \mathcal{L}_k(W(\Lambda)) &\approx \mathcal{L}_k(W_k) \\
    &+ \nabla \mathcal{L}_k(W_k)^T (W(\Lambda) - W_k) \\
    &+ \frac{1}{2} (W(\Lambda) - W_k)^T H_k (W(\Lambda) - W_k)
\end{split}
\end{equation}
where $H_k = \nabla^2 \mathcal{L}_k(W_k) \in \mathbb{R}^{D \times D}$ is the full, non-diagonal Hessian matrix of task $k$ evaluated at $W_k$.

To make this optimization tractable, we introduce two key approximations:
\begin{assumption}[Layer Block-Diagonal Hessian]
    We assume that cross-layer second-order interactions are negligible. That is, the Hessian $H_k$ is block-diagonal across the $L$ layers:
    \begin{equation}
        H_k \approx \text{diag}\left(H_k^1, H_k^2, \dots, H_k^L\right)
    \end{equation}
    where $H_k^l \in \mathbb{R}^{d_l \times d_l}$ is the layer-specific Hessian for layer $l$ of dimension $d_l$.
\end{assumption}

Under Assumption 3.1, the Taylor expansion decomposes across layers:
\begin{equation}
\begin{split}
    \mathcal{L}_k(W(\Lambda)) &\approx \mathcal{L}_k(W_k) + \sum_{l=1}^L \Big[ (g_{k,0}^l)^T (W^l(\Lambda^l) - W_k^l) \\
    &+ \frac{1}{2} \left(W^l(\Lambda^l) - W_k^l\right)^T H_k^l \left(W^l(\Lambda^l) - W_k^l\right) \Big]
\end{split}
\end{equation}
where $g_{k,0}^l = \nabla_{W^l} \mathcal{L}_k(W_k) \in \mathbb{R}^{d_l}$ is the unperturbed residual gradient of task $k$ at layer $l$.

Our second, and most critical, technique is the \textbf{Low-Dimensional Subspace Projection}. Rather than searching the massive $d_l$-dimensional layer parameter space, we restrict our search space strictly to the $K$-dimensional subspace spanned by the task vectors $v_1^l, \dots, v_K^l$. Let $V^l = [v_1^l, v_2^l, \dots, v_K^l] \in \mathbb{R}^{d_l \times K}$ be the matrix of task vectors. We can express the merged parameters at layer $l$ as $W^l(\Lambda^l) = W_0^l + V^l \Lambda^l$. The parameter shift from expert $k$ at layer $l$ is then:
\begin{equation}
\begin{split}
    W^l(\Lambda^l) - W_k^l &= \left(W_0^l + V^l \Lambda^l\right) - \left(W_0^l + v_k^l\right) \\
    &= V^l \Lambda^l - v_k^l
\end{split}
\end{equation}

Substituting this back into the joint loss optimization problem, we see that the joint multi-task loss minimization objective decomposes into $L$ independent layer-wise quadratic minimization problems:
\begin{equation}
\begin{split}
    \min_{\Lambda^l} \sum_{k=1}^K \alpha_k \Big[ &(g_{k,0}^l)^T (V^l \Lambda^l - v_k^l) \\
    + &\frac{1}{2} \left(V^l \Lambda^l - v_k^l\right)^T H_k^l \left(V^l \Lambda^l - v_k^l\right) \Big]
\end{split}
\end{equation}

\subsection{Derivation of the Analytical Optimal Solution}
We can expand the layer-wise quadratic objective function $f(\Lambda^l)$:
\begin{equation}
\begin{split}
    f(\Lambda^l) &= \sum_{k=1}^K \alpha_k \Big[ (\Lambda^l)^T (V^l)^T g_{k,0}^l \\
    &+ \frac{1}{2} (\Lambda^l)^T (V^l)^T H_k^l V^l \Lambda^l - (\Lambda^l)^T (V^l)^T H_k^l v_k^l \Big] \\
    &+ C'
\end{split}
\end{equation}
where $C'$ is a constant independent of $\Lambda^l$. Isolating the terms dependent on $\Lambda^l$, we can write:
\begin{equation}
    f(\Lambda^l) = \frac{1}{2} (\Lambda^l)^T A^l \Lambda^l - (\Lambda^l)^T (b^l - d^l) + C
\end{equation}
where $C$ is a constant, and the projected Hessian matrix $A^l \in \mathbb{R}^{K \times K}$, projected vector $b^l \in \mathbb{R}^K$, and projected first-order gradient vector $d^l \in \mathbb{R}^K$ are defined as:
\begin{align}
    A^l &= \sum_{k=1}^K \alpha_k (V^l)^T H_k^l V^l \\
    b^l &= \sum_{k=1}^K \alpha_k (V^l)^T H_k^l v_k^l \\
    d^l &= \sum_{k=1}^K \alpha_k (V^l)^T g_{k,0}^l
\end{align}

Specifically, the individual elements are computed as:
\begin{align}
    A^l_{ij} &= \sum_{k=1}^K \alpha_k (v_i^l)^T H_k^l v_j^l \quad \text{for } i, j \in \{1, \dots, K\} \label{eq:a_matrix} \\
    b^l_i &= \sum_{k=1}^K \alpha_k (v_i^l)^T H_k^l v_k^l \quad \text{for } i \in \{1, \dots, K\} \label{eq:b_vector} \\
    d^l_i &= \sum_{k=1}^K \alpha_k (v_i^l)^T g_{k,0}^l \quad \text{for } i \in \{1, \dots, K\} \label{eq:d_vector}
\end{align}

To ensure numerical stability and invertibility of $A^l$ even in low-curvature or redundant subspaces, we incorporate Ridge (Tikhonov) regularization with a small penalty parameter $\gamma > 0$:
\begin{equation}
\begin{split}
    f_{\text{reg}}(\Lambda^l) &= \frac{1}{2} (\Lambda^l)^T A^l \Lambda^l - (\Lambda^l)^T (b^l - d^l) \\
    &+ \frac{\gamma}{2} \|\Lambda^l\|_2^2
\end{split}
\end{equation}
Taking the gradient of $f_{\text{reg}}(\Lambda^l)$ with respect to $\Lambda^l$ and setting it to zero:
\begin{equation}
\begin{split}
    &\nabla_{\Lambda^l} f_{\text{reg}}(\Lambda^l) = A^l \Lambda^l - (b^l - d^l) + \gamma \Lambda^l = 0 \\
    &\implies (A^l + \gamma I) \Lambda^l = b^l - d^l
\end{split}
\end{equation}
This yields our final, exact, **closed-form optimal solution** under non-zero expert gradients:
\begin{equation}
    \Lambda^{l, *} = (A^l + \gamma I)^{-1} (b^l - d^l)
\end{equation}

\subsection{Scale-Normalized ACM (ACM-Norm)}
The raw analytical system defined above accumulates Hessian contributions directly across tasks. Because different tasks can have vastly different loss scales and gradient/Hessian magnitudes, the joint system is often heavily dominated by the task with the largest loss or gradient scale (e.g., MNIST). This scale imbalance causes a severe \emph{sacrificial task bias}, where the solved coefficients are highly optimized for the dominant task while degrading others.

To achieve scale-invariance and ensure a balanced multi-task merge, we propose \textbf{Scale-Normalized ACM (ACM-Norm)}. In this variant, we normalize each task's projected Hessian contribution by its trace before joint assembly. Let $A_k^l \in \mathbb{R}^{K \times K}$ represent the projected Hessian of task $k$ at layer $l$:
\begin{equation}
    (A_k^l)_{ij} = (v_i^l)^T H_k^l v_j^l
\end{equation}
We define the scale-normalized system matrix $\tilde{A}^l$, target vector $\tilde{b}^l$, and first-order linear term $\tilde{d}^l$ as:
\begin{align}
    \tilde{A}^l &= \sum_{k=1}^K \alpha_k \frac{A_k^l}{\text{Tr}(A_k^l)} \\
    \tilde{b}^l_i &= \sum_{k=1}^K \alpha_k \frac{(A_k^l)_{ik}}{\text{Tr}(A_k^l)} \\
    \tilde{d}^l_i &= \sum_{k=1}^K \alpha_k \frac{(v_i^l)^T g_{k,0}^l}{\text{Tr}(A_k^l)}
\end{align}
where $\text{Tr}(A_k^l) = \sum_{j=1}^K (v_j^l)^T H_k^l v_j^l$ is the trace of the projected Hessian, representing the total second-order parameter sensitivity of task $k$ at layer $l$. This scale-invariant system is then solved under Ridge regularization:
\begin{equation}
    \Lambda^{l, *}_{\text{Norm}} = (\tilde{A}^l + \gamma I)^{-1} (\tilde{b}^l - \tilde{d}^l)
\end{equation}
This advanced formulation guarantees that no single task dominates the fusion process, ensuring balanced multi-task performance.

\subsection{Global-Normalized ACM (ACM-GlobalNorm)}
While Scale-Normalized ACM (ACM-Norm) successfully resolves loss scale imbalance by trace-normalizing the projected Hessian layer-by-layer, it introduces a physical distortion: it forces the total sensitivity (trace of the normalized system) of every layer to be exactly $1.0$ for all tasks. This flattens and destroys the relative layer-wise parameter sensitivity profile of each individual task (i.e., different layers naturally have vastly different sensitivities to weight changes, and this crucial structure is lost).

To achieve scale-invariance across tasks while fully preserving the relative layer sensitivities of different layers within each task, we propose \textbf{Global-Normalized ACM (ACM-GlobalNorm)}. In this variant, we normalize each task's projected Hessian contribution by its \emph{global trace} (sum of traces across all layers) rather than normalizing layer-by-layer. Let $T_k = \sum_{l'=1}^L \text{Tr}(A_k^{l'})$ represent the global projected sensitivity of task $k$. We define the globally scale-normalized system matrix $\bar{A}^l$, target vector $\bar{b}^l$, and first-order linear term $\bar{d}^l$ at layer $l$ as:
\begin{align}
    \bar{A}^l &= \sum_{k=1}^K \alpha_k \frac{A_k^l}{T_k} \\
    \bar{b}^l_i &= \sum_{k=1}^K \alpha_k \frac{(A_k^l)_{ik}}{T_k} \\
    \bar{d}^l_i &= \sum_{k=1}^K \alpha_k \frac{(v_i^l)^T g_{k,0}^l}{T_k}
\end{align}
The optimal globally scale-normalized coefficients are then solved under Ridge regularization:
\begin{equation}
    \Lambda^{l, *}_{\text{GlobalNorm}} = (\bar{A}^l + \gamma I)^{-1} (\bar{b}^l - \bar{d}^l)
\end{equation}
This formulation preserves the relative depth-wise sensitivity profile of each expert while establishing joint scale-variance correction, successfully balancing multi-task fusion without sacrificial bias or layer-wise sensitivity flattening.

\begin{theorem}
    If the regularized projected Hessian $(A^l + \gamma I)$ is positive definite (which is guaranteed for any $\gamma > 0$ as $A^l$ is positive semi-definite), the layer-wise coefficient vector $\Lambda^{l, *}$ is the unique global minimizer of the regularized projected joint quadratic loss landscape.
\end{theorem}
\begin{proof}
    The Hessian of $f_{\text{reg}}(\Lambda^l)$ is $\nabla^2_{\Lambda^l} f_{\text{reg}}(\Lambda^l) = A^l + \gamma I$. Since $A^l = \sum_k \alpha_k (V^l)^T H_k^l V^l$, and the Hessian $H_k^l$ is positive semi-definite near its local minimum, the projected matrix $A^l$ is positive semi-definite. Adding $\gamma I$ with $\gamma > 0$ ensures that all eigenvalues of $A^l + \gamma I$ are strictly positive (bounded below by $\gamma$). Thus, the regularized objective is strictly convex, and setting the gradient to zero yields the unique global minimum.
\end{proof}

\subsection{Finite-Difference Projected Hessian-Vector Estimation}
Directly evaluating Equations \ref{eq:a_matrix} and \ref{eq:b_vector} is impossible because we cannot construct or store the full $D \times D$ Hessian matrices $H_k$. To bypass this, we utilize a cheap finite-difference scheme. Let $\epsilon > 0$ be a small perturbation scale. Using a first-order Taylor expansion of the task gradient $\nabla \mathcal{L}_k$ around the expert weights $W_k$:
\begin{equation}
    \nabla \mathcal{L}_k(W_k + \epsilon v_j) \approx \nabla \mathcal{L}_k(W_k) + \epsilon H_k v_j
\end{equation}

In standard settings, assuming the expert weights $W_k$ have reached absolute local convergence, the unperturbed gradient vanishes: $\nabla \mathcal{L}_k(W_k) \approx 0$. However, in practice, neural network training rarely reaches a perfect zero-gradient minimum. If a non-zero residual gradient $g_{k,0} = \nabla \mathcal{L}_k(W_k) \neq 0$ is present, directly evaluating $H_k v_j \approx \frac{1}{\epsilon} \nabla \mathcal{L}_k(W_k + \epsilon v_j)$ results in severe numerical instability, where the optimization noise is amplified by $1/\epsilon$.

To resolve this instability and ensure complete robustness to partially converged expert checkpoints, we propose \textbf{Gradient Subtraction ACM}. By explicitly computing the unperturbed expert gradient $g_{k,0} = \nabla \mathcal{L}_k(W_k)$, we subtract it from the perturbed gradient to cancel out the residual gradient term:
\begin{equation}
    H_k v_j \approx \frac{1}{\epsilon} \left[ \nabla \mathcal{L}_k(W_k + \epsilon v_j) - \nabla \mathcal{L}_k(W_k) \right]
\end{equation}
Thus, we can estimate the projected Hessian-vector product $(v_i^l)^T H_k^l v_j^l$ layer-by-layer as:
\begin{equation}
    (v_i^l)^T H_k^l v_j^l \approx \frac{1}{\epsilon} \langle v_i^l, g_{k,j}^l - g_{k,0}^l \rangle
\end{equation}
where $g_{k,j}^l = \nabla_{W^l} \mathcal{L}_k(W_k + \epsilon v_j) \in \mathbb{R}^{d_l}$ is the perturbed gradient, and $g_{k,0}^l = \nabla_{W^l} \mathcal{L}_k(W_k) \in \mathbb{R}^{d_l}$ is the unperturbed gradient of the task-specific loss $\mathcal{L}_k$ with respect to the weights of layer $l$. This subtraction mathematically cancels the residual gradient and bounds the truncation error to a clean $O(\epsilon)$ (see Appendix B for the formal proof).

This finite-difference scheme requires only $K^2 + K$ gradient computations (one for each task $k$ and direction $j$, plus $K$ unperturbed gradient computations). During calibration, we feed a single small calibration batch (e.g., 32 samples) of task $k$'s dataset to the perturbed model, perform a forward-backward pass, and accumulate the scalar dot products layer-by-layer. The entire calibration and system solving process takes less than 5 seconds on a single GPU.

\subsection{Computational and Space Complexity Analysis}
The computational and space complexities of ACM are highly favorable:
\begin{itemize}
    \item \textbf{Space Complexity:} For each layer $l$, we store a $K \times K$ matrix $A^l$ and a $K$-dimensional vector $b^l$. The overall storage cost is $O(L K^2)$ scalars, which is completely negligible (e.g., for $L=14, K=4$, this is only 280 floating-point numbers).
    \item \textbf{Inversion Complexity:} Solving the linear system requires inverting the $K \times K$ matrix $A^l + \gamma I$, which has a computational cost of $O(K^3)$ per layer. Since $K \ll D$, this inversion is practically instantaneous.
    \item \textbf{Calibration Complexity:} The finite-difference gradient computation requires $K^2$ backward passes on a tiny calibration batch of size 32. This is extremely cheap compared to Test-Time Adaptation methods like AdaMerging, which require $500$ to $1000$ backward passes over the entire test dataset.
\end{itemize}
This complete formulation demonstrates the theoretical and operational elegance of ACM: it is a training-free, computationally trivial, and mathematically exact solution to multi-task parameter fusion.
